

\section{Conclusion}
%We suggest a novel approach: grammar-constrained decoding to improve the few-shot performance of LLMs on structured NLP tasks where they struggle.
%We suggest a novel approach, grammar-constrained decoding, to improve the few-shot performance of LLMs on structured NLP tasks, which pose challenges to them.
This work introduces grammar-constrained decoding (GCD) for enhancing the few-shot performance of LLMs on challenging structured NLP tasks.
%We show that many NLP tasks can be formulated as formal grammars and GCD can be used to improve the performance of LLMs on these tasks.
We showed that many NLP tasks can be formulated as formal grammars, and that GCD can enhance the performance of LLMs on these tasks.
%With input-dependent grammar, we further extend the scope of GCD to tasks that require different output structures for different inputs.
With input-dependent grammars, we further broaden the scope of GCD to accommodate tasks where the set of valid output structures is constrained by the given input.
%Our experiments show that while LLMs struggle with tasks that require structured output, GCD can significantly improve the performance of LLMs on these tasks.
Our experiments indicate that, whereas unconstrained LLMs have difficulty tackling tasks requiring structured outputs, GCD can significantly bolster the performance of LLMs on such tasks.
%We view GCD as a rapid and cost-effective adaptation strategy that enables LLMs to produce reliable structured outputs without the need for finetuning.
We envision GCD as a swift and cost-efficient adaptation strategy, allowing LLMs to generate reliable structured outputs without the necessity for costly and cumbersome finetuning.
%Given the rapid advancements in LLMs, we believe that GCD will emerge as a potent tool for dealing with structured NLP tasks.
Given the fast pace of LLM evolution, we anticipate the usefulness of GCD for boosting pretrained LLMs to further increase over time.

\xhdr{Best practices for GCD}
We conclude with considerations regarding the effective use of GCD.

\begin{enumerate}
    \denselist
    \item
    % It is advisable to use the best available LLMs. We see that
    GCD is more effective with larger LLMs. When possible, use the largest available LLM.
    \item Grammars should be as restrictive as possible. Consider using input\hyp dependent grammars.
    \item While GCD is broadly applicable to many tasks, it is not a silver bullet. Tasks that require syntactic understanding of the input (e.g., constituency parsing) are less suitable for GCD.
    \item In the presence of task\hyp specific training data, finetuning a small model may still yield better performance, at the cost of decreased convenience (\cf\ \Appref{app_sec:finetuning_vs_gcd}).
\end{enumerate}

